\section{Recognition of the intrinsic period lattice}
\label{sec:periods}
The scalar inverse recovers a physical patch before periodicity is used.
We next give the local period argument needed to extend that patch.
This argument is retained from the preceding repeated-motif revision
\cite{Supplement}; it is included here so that the primary theorem does
not depend on a different sensor's acquisition proof.

\begin{definition}[Bounded periodic class]\label{def:class}
Fix $r_0,L_0,V_0,D_0,d_0,\kappa_-,\kappa_+>0$. The class
$\mathcal B$ consists of nonempty planar unions
\begin{equation}\label{eq:period-presentation}
 \mathcal O=\bigcup_{i=1}^r(C_i+\Lambda_0),\qquad
 \Lambda_0=L\Z^2,\qquad 1\le r\le r_0,
\end{equation}
with exactly $r$ component orbits modulo $\Lambda_0$,
$\|L\|,\|L^{-1}\|\le L_0$, and $\covol\Lambda_0\ge V_0$.
Every physical component is a compact strictly convex body of diameter
at most $D_0$, with curvature in $[\kappa_-,\kappa_+]$; distinct
components are separated by at least $d_0$. Centered support functions
have $C^{6,\beta}$ norm at most $M_6$, for a fixed $0<\beta\le1$.
The unknown presentation need not be primitive. Repeated shapes and
arbitrary symmetries are allowed. Its full period group is
$\Pi(\mathcal O)=\{w:\mathcal O+w=\mathcal O\}$.
\end{definition}

The high-order bound will only enter finite $C^2$ reconstruction.
Set $B=1+L_0$ and $Q_s=[-s,s]^2$. A convenient fixed observation and
launch aperture is
\begin{equation}\label{eq:aperture}
 R_0=12B+3D_0+4,\qquad R=R_0+mt+2,\qquad Q_R,
\end{equation}
where $t,m$ are from \eqref{eq:design}. Scalar preparations centered
on $Q_{R_0}$ and on the forward chains needed for it are supported in
$Q_R$ for every $\sigma\le\sigma_0$. Each physical flight has the same
fixed duration $t$. There is no recentering at an estimated obstacle.

For a component $C$ write $z_C$ for its Steiner point and $p_C$ for its
centered support function. In the laboratory directions define
\begin{align}
 d_*(C,D)&=\|z_C-z_D\|_\infty+\|p_C-p_D\|_\infty,\label{eq:bodymetric}\\
 \mathcal D_{\mathcal O}(w)&=
 \max_{C:z_C\in Q_B}\ \max_{s\in\{-1,1\}}
                   \inf_{D\in\mathcal C}d_*(C+sw,D),\label{eq:defect}
\end{align}
where $\mathcal C$ is the set of physical components. The metric does
not optimize over rotations. Hausdorff error $e$ implies center error
at most $2e$ and centered-support error at most $3e$. All relevant
infima are attained, since distant centers cannot minimize them.

\begin{lemma}[A finite test for a global period]\label{lem:period}
One has $\mathcal D_{\mathcal O}(w)=0$ if and only if
$w\in\Pi(\mathcal O)$. For $\|w\|_\infty\le8B$ the zero test uses
only components with centers in $Q_B$ and $Q_{9B}$.
\end{lemma}
\begin{proof}
Every $\Lambda_0$-component orbit has a representative with center in
$L[-1/2,1/2]^2\subset Q_B$. If the defect vanishes, each such
representative has both translates by $+w$ and $-w$ among the physical
components. Translating by $\Lambda_0$ gives
$\mathcal O+w\subset\mathcal O$ and $\mathcal O-w\subset\mathcal O$.
The latter implies the reverse of the former inclusion. The converse
is immediate. For the bounded zero test, every possible equal target
has center $z_C\pm w\in Q_{9B}$.
\end{proof}

\begin{lemma}[A protected generating list]\label{lem:generators}
The group $\Pi=\Pi(\mathcal O)$ is a lattice containing $\Lambda_0$,
\begin{equation}\label{eq:index}
 [\Pi:\Lambda_0]\le r\le r_0,\qquad \covol\Pi\ge V_*:=V_0/r_0.
\end{equation}
For any root component centered in $Q_{2B}$, take its center differences
to components centered in $Q_{4B}$ and retain those passing
Lemma~\ref{lem:period}. Their integer span is exactly $\Pi$.
\end{lemma}
\begin{proof}
Map a period coset modulo $\Lambda_0$ to the $\Lambda_0$-orbit of the
translated root. This is injective, since a compact body has no nonzero
translation symmetry. There are at most $r$ such cosets. A finite
extension of a full-rank lattice is a full-rank lattice; the covolume
formula proves \eqref{eq:index}. This quotient group acts freely on
the $r$ component orbits, so its order divides $r$.

All accepted differences are true periods. Conversely, the two columns
of $L$ and representatives of all cosets of $\Pi/\Lambda_0$ can be
chosen to have norm less than $B$, using the same centered fundamental
parallelogram. Translating the root by any of these vectors gives a
target in $Q_{3B}$, strictly inside the cutoff $Q_{4B}$. The accepted
list contains generators of $\Lambda_0$ and representatives of every
coset, and hence generates $\Pi$.
\end{proof}

\begin{theorem}[Whole-table rigidity from scalar collision laws]
\label{thm:global-exact}
On $\mathcal B$, the collision-probability functionals for the two
opposed directions at the fixed duration \eqref{eq:design}, with
controlled smooth starting densities in $Q_R$, determine the entire
obstacle union, its full period lattice $\Pi$, its primitive number
$r_*$ of component orbits, and primitive free area $A_*$. The same is
true for the family of localized kernel commands at all sufficiently
small scales. No impact position, impact angle, collision time, shape
label, integer copy label, primitive cell or free-area normalization
is an observation. The laboratory frame is supplied by the launch
controls; forgetting it leaves a common Euclidean-isometry ambiguity.
\end{theorem}
\begin{proof}
Corollary~\ref{cor:local} recovers the occupation indicator on
$Q_{R_0}$, and regular closedness identifies the physical set there.
Every component with center in $Q_{10B}$, and every component needed
to recognize one of these, lies wholly within the protective collar
of this square. Physical separation identifies the complete components.
Lemmas~\ref{lem:period} and \ref{lem:generators} then recover $\Pi$.

All $\Pi$-component orbits have representatives in $Q_B$. Among
components centered in $Q_{2B}$, two are in the same orbit exactly when
their center difference passes the period test. Indeed a period maps
the first to a physical component with the second's Steiner point;
two disjoint full-dimensional convex bodies cannot have the same
Steiner point, which lies in their interiors. Thus the relation gives
$r_*$ and one representative per orbit. These and $\Pi$ determine the
union everywhere. Finally
\begin{equation}\label{eq:area}
 A_*=\covol\Pi-\sum_{j=1}^{r_*}\area(C_j).
\end{equation}
The conclusion concerns the intrinsic union, not a chosen multiple-cell
presentation. Basis computation from finite exact period vectors is
justified quantitatively in Lemma~\ref{lem:locking} below.
\end{proof}

\begin{definition}[Nonperiod patch margin]\label{def:margin}
For $\eta>0$, let $\mathcal B_\eta$ consist of those members of
$\mathcal B$ such that
\begin{equation}\label{eq:margin}
 w=z_D-z_C,\quad z_C\in Q_{3B},\ z_D\in Q_{5B},\quad
 w\notin\Pi(\mathcal O)\quad\Longrightarrow\quad
                  \mathcal D_{\mathcal O}(w)\ge\eta.
\end{equation}
This concerns translations of a whole central patch, not separation
between individual body shapes. For each fixed table it holds with some
positive $\eta$, because there are finitely many candidate pairs and
Lemma~\ref{lem:period} makes every nonperiod defect positive.
\end{definition}

The next deterministic statement is the link between geometric estimation
and exact discrete information. It does not suppose that perturbed real
vectors themselves generate a discrete group.
\begin{lemma}[Robust period classification and rational locking]
\label{lem:locking}
Fix the numerical bounds of $\mathcal B_\eta$. Suppose complete
components on the protected patch are given by matched approximations
of Hausdorff error at most $e$. For sufficiently small $e$, depending
on $\eta$ and the geometric bounds, finite comparisons recover the true
primitive orbit count and the exact rational coordinates of an accepted
period-generating list in an independent-pair basis. They give a matched
basis for $\Pi$ with error $O(e)$ and an area estimate with error $O(e)$.
\end{lemma}
\begin{proof}
Use approximate centers in $Q_{2B}$ as sources, choose a root there,
and use its targets in $Q_{4B}$. A measured difference $\widehat w$
has a corresponding true center difference $w$ with error at most
$4e$. Test both signs of this translation on all selected sources,
minimizing $d_*$ over approximate components centered in $Q_{10B}$.
For each source-target comparison, the center term changes by at most
$8e$ and the centered-support term by at most $6e$.

If $w$ is a period, all needed partners lie within the target cutoff
by the strict aperture margins, and the measured defect is at most
$14e$. If it is not, its true endpoints lie in $Q_{3B},Q_{5B}$, and
\eqref{eq:margin} applies. Every witness with true center in $Q_B$ is
selected. Restricting the targets cannot reduce a true positive infimum,
so the measured defect is at least $\eta-14e$. Reserve another $6e$
for numerical enclosures of centers and support suprema. Threshold
$\eta/2$, with $20e<\eta/4$, accepts exactly the true periods among
these candidates. The same comparison on all pairs of selected sources
gives their true orbit equivalence relation and hence $r_*$. Targets
for these tests also lie within $Q_{10B}$. Lemma~\ref{lem:generators}
and its strict cutoff margin ensure all protected generators are present.

All accepted vectors have norm at most $U=12B$ for small $e$. A nonzero
determinant of two true periods is an integer multiple of
$\covol\Pi\ge V_*$. Determinant perturbations are $O(Ue)$; after
shrinking $e$, independence is decided at threshold $V_*/2$. Choose
an independent accepted pair with true matrix $D$. The sublattice
$D\Z^2$ has index
\[
       n=|\det D|/\covol\Pi\le Q:=\max(1,\lceil U^2/V_*\rceil).
\]
Every coordinate of $D^{-1}w_j$ has reduced denominator dividing $n$,
and a prior-bounded numerator. The inverse-matrix identity bounds its
estimated error by $C_*e$. Distinct rationals of denominator at most
$Q$ differ by at least $Q^{-2}$. If $C_*e<1/(3Q^2)$, finite search
therefore identifies each exact coordinate uniquely.

Let $q_*$ be their common denominator; it divides $n$ and is at most
$Q$. A column Hermite basis $H$ of the integer span of $q_*I_2$ and
$q_*D^{-1}w_j$ gives
\begin{equation}\label{eq:hermite}
                  \Pi=(DH/q_*)\Z^2.
\end{equation}
The integer group contains $q_*\Z^2$, so the Hermite entries are
bounded in terms of $Q$. Replacing $D$ by its measured version gives
a matched basis with error $O(e)$. This basis need not be a continuously
chosen reduced basis. The parallel-body bound gives an $O(e)$ area
error for each matched convex component; use \eqref{eq:area} for
$A_*$. Certified finite angular meshes and the uniform support Lipschitz
bound enclose all suprema and center integrals to the reserved precision.
\end{proof}
